\subsection{A shared component in every scoring output}
\label{sec:axis}

We computed eight scoring outputs on one identical block of retained DOCKSTRING poses,
2{,}000 ligands by nine targets, with the top AutoDock Vina pose held fixed in each of the
18{,}000 cells. Pose selection is therefore common to all eight, and whatever separates them
comes from the scoring function alone.

Across the eight, the leading principal component of the target correlation matrix carries
between 0.414 and 0.669 of the variance. Its loading vector has cosine 0.950 to 0.986
with the uniform target direction, and no loading is negative in any output. Thus, in every scorer
tested, the largest single component is a common rise and fall across targets. Removing each
ligand's mean across targets raises the participation-ratio dimension
in all eight (Figure~\ref{fig:axis}).

Four descriptive criteria were written down before the non-Vina scores were inspected: enough
variance in PC1, a nearly uniform positive loading vector, a low raw participation ratio, and a
clear rise after centring. Seven outputs meet all four on their point estimates. PLECscore misses
the PC1 and raw-dimension gates, although only the PC1 miss is clear relative to its uncertainty.
The complete criteria and intervals are reported in Supplementary Section S3.

The size of the shared component varies with scorer feature class (Figure~\ref{fig:axis}a).
Atom-pair contact-count models are highest, empirical interaction-term sums follow, and NNScore
and PLECscore are lower. This ordering is consistent with contact counts loading more strongly
on molecular bulk, but the related model lineages and common retained poses prevent a causal
interpretation.

The same ordering holds under every leave-one-target-out analysis and most larger target subsets,
but weakens on the smallest subsets. The eight outputs are not eight independent functions:
Vina and smina evaluate the same function, Vinardo reweights the same empirical form, the
three RF-Scores are versions of one model, and RF-Score v3 takes the Vina intermolecular
terms as input features, which leaves about four scoring lineages.

The same pattern appears in the two full Vina matrices, each with its own target count. On
Docking-44 (12{,}651 ligands, 44 targets) PC1 carries 0.733 of the target-correlation
variance at cosine 0.970 to the uniform direction, and the raw participation ratio is 1.834.
On DOCKSTRING-58 (260{,}060 ligands, 58 targets) the three values are 0.659, 0.969 and 2.266.
A 44-column matrix, in other words, behaves like fewer than two independent columns.

\subsection{What remains after the offset is removed}

Removing each ligand's mean across targets leaves a residual with participation-ratio
dimension 9.295 on Docking-44 and 18.158 on DOCKSTRING-58, each measured on the same
deterministic 12{,}000-ligand support used for that matrix's nulls. Part of this rise is
mechanical: the operation caps the residual dimension at $P-1$, which is 43
and 57 for these two matrices, and it pushes any estimate upward. The question is whether a
matrix with no target-resolved structure at all would also land near 9 or 18 after the same
steps. The matched nulls do not.

The Docking-44 complete-case value is higher. Supplementary Section S2 shows that this reflects
a chemically different ligand support, driven largely by one incomplete target column, rather
than an imputation-only effect.

We ran four transformation-matched null constructions through the same pipeline on the
same supports, 500 repetitions each (Figure~\ref{fig:residual}a). An additive null adds
independent Gaussian noise with the observed target-specific residual standard deviations to
fitted grand, ligand and target effects. A column-permutation null permutes each observed
residual target column independently across ligands. A direction null replaces each ligand's
residual by a random direction in the zero-sum subspace, with its norm preserved. A
size-conditioned null permutes each processed raw target column within ten equally populated
molecular-weight bins, which keeps each target's score distribution inside a coarse size
domain while destroying any cross-target correspondence for a given ligand.

The three nulls that ignore molecular size sit against the ceiling. Their median residual
dimensions are 42.400, 42.403 and 42.628 on Docking-44, and 56.482, 56.481 and 56.612 on
DOCKSTRING-58, with 95\% simulation ranges at most 0.11 wide. Both observed values fall below
all 500 repetitions of all three. The size-conditioned null is the only construction that
moves away from the ceiling, to medians of 34.875 and 47.039, which are 0.81 and 0.83 of the
bound. Coarse ligand size therefore accounts for part of the concentration, but not most of it:
the size-conditioned null closes 23.8\% of the distance between the one-bin version and
the observed value on Docking-44, and 25.1\% on DOCKSTRING-58.

The residual map does depend on the library it was estimated from. Residual maps estimated on
the lower and upper molecular-weight quartiles of the same library agree at rank correlation
0.205 for Docking-44 (all 12{,}651 ligands) and 0.351 for DOCKSTRING-58 (a 15{,}000-ligand
random support), shown in Figure~\ref{fig:residual}b. Across 200 size-matched splits of the
same library into halves sharing no whole chemical group, the same comparison gives 0.992 and
0.994. These controls localise the disagreement to the sampled molecular-weight domain, but
molecular weight was selected post hoc and covaries with other chemical properties; it is not
identified as the cause. Between the two quartile maps, 42.5\% of target-pair correlations
change sign on Docking-44 and 36.7\% on DOCKSTRING-58. Estimate the map on the library that
will be screened.
